\documentclass[a4paper,12pt]{article}

\usepackage[utf8]{inputenc}
\usepackage[T5]{fontenc}
\usepackage[vietnamese]{babel}
\usepackage{amsmath, amssymb}
\usepackage{geometry}
\usepackage{enumitem}
\usepackage{tcolorbox}

\geometry{margin=2cm}

\begin{document}

\begin{center}
    \textbf{\Large DUYỆT ĐỒ THỊ THEO BFS}
\end{center}

\section*{ĐỀ BÀI}
Bạn có một mạng lưới các thành phố được nối với nhau bằng các con đường hai chiều.

Bạn muốn bắt đầu từ một thành phố và khám phá tất cả các thành phố có thể đi tới theo cách mở rộng dần (BFS).

Cho đồ thị vô hướng gồm $n$ đỉnh và $m$ cạnh. Hãy in ra thứ tự duyệt BFS bắt đầu từ đỉnh $s$.

\section*{DỮ LIỆU VÀO}
\begin{itemize}
    \item Ba số nguyên $n, m, s$ $(1 \le n \le 10^5,\ 1 \le m \le 2 \times 10^5)$.
    \item $m$ dòng tiếp theo, mỗi dòng gồm hai số nguyên $u, v$ biểu diễn một cạnh giữa $u$ và $v$.
\end{itemize}

\section*{DỮ LIỆU RA}
\begin{itemize}
    \item In ra thứ tự các đỉnh được duyệt theo BFS bắt đầu từ đỉnh $s$.
\end{itemize}

\section*{SAMPLE INPUT}
\begin{tcolorbox}
5 5 1 \\
1 2 \\
1 3 \\
2 4 \\
3 5 \\
4 5
\end{tcolorbox}

\section*{SAMPLE OUTPUT}
\begin{tcolorbox}
1 2 3 4 5
\end{tcolorbox}

\section*{SUBTASK}
\begin{itemize}
    \item 50\%: $n \le 1000$
    \item 100\%: $n \le 10^5$
\end{itemize}
\newpage
\section*{HƯỚNG DẪN GIẢI}
Đây là bài toán duyệt đồ thị theo chiều rộng (BFS).

\textbf{Ý tưởng:}
\begin{itemize}
    \item Sử dụng một hàng đợi (queue) để lưu các đỉnh cần duyệt.
    \item Bắt đầu từ đỉnh $s$, đưa $s$ vào hàng đợi và đánh dấu đã thăm.
    \item Lặp lại cho đến khi hàng đợi rỗng:
    \begin{itemize}
        \item Lấy đỉnh ở đầu hàng đợi ra và in ra.
        \item Xét tất cả các đỉnh kề của nó.
        \item Nếu đỉnh kề chưa được thăm, đánh dấu và đưa vào hàng đợi.
    \end{itemize}
\end{itemize}

\textbf{Lưu ý:}
\begin{itemize}
    \item Nên đánh dấu đã thăm ngay khi đưa vào hàng đợi để tránh lặp.
    \item Có thể sắp xếp danh sách kề để đảm bảo thứ tự duyệt tăng dần.
\end{itemize}

\textbf{Độ phức tạp:} $O(n + m)$

\end{document}